\DocumentMetadata{
  pdfversion=2.0,
  pdfstandard=ua-2,
  lang=en-US,
  tagging=on,
  tagging-setup={math/setup=mathml-SE,tabsorder=structure}
}
\documentclass[11pt,letterpaper]{article}
\usepackage[margin=0.68in,includefoot]{geometry}
\usepackage{newcomputermodern}
\usepackage{amsmath,amscd,mathtools}
\providecommand{\square}{\mdlgwhtsquare}
\usepackage[hidelinks]{hyperref}
\hypersetup{
  pdftitle={Analysis First-Year Exam, January 2015, Part 1},
  pdfauthor={University of Florida Department of Mathematics},
  pdfsubject={Graduate mathematics examination},
  pdfdisplaydoctitle=true,
  bookmarksopen=true
}
\setcounter{secnumdepth}{0}
\setlength{\parindent}{0pt}
\setlength{\parskip}{0.28em}
\setlength{\emergencystretch}{3em}
\setlength{\leftmargini}{2.15em}
\setlength{\leftmarginii}{2.65em}
\setlength{\leftmarginiii}{3em}
\renewcommand{\labelenumi}{\textbf{\theenumi.}}
\pagestyle{empty}
\raggedbottom
\allowdisplaybreaks
\newenvironment{examproblems}[1]{%
  \begin{enumerate}%
  \setcounter{enumi}{#1}%
  \setlength{\topsep}{0.35em}%
  \setlength{\partopsep}{0pt}%
  \setlength{\itemsep}{0.48em}%
  \setlength{\parsep}{0.18em}%
}{\end{enumerate}}
\newenvironment{examsubparts}{%
  \begin{enumerate}%
  \setlength{\topsep}{0.18em}%
  \setlength{\partopsep}{0pt}%
  \setlength{\itemsep}{0.16em}%
  \setlength{\parsep}{0.1em}%
}{\end{enumerate}}
\newenvironment{examinstructions}{%
  \par\begingroup\itshape
}{%
  \par\endgroup\vspace{0.18em}
}
\makeatletter
\renewcommand\section{\@startsection{section}{1}{0pt}%
  {-1.25ex plus -.25ex minus -.1ex}%
  {0.55ex plus .1ex}%
  {\normalfont\large\bfseries}}
\renewcommand{\@maketitle}{%
  \newpage\null\vskip -1.2em%
  {\footnotesize\bfseries
    \href{https://www.ufl.edu/}{University of Florida}, \href{https://math.ufl.edu/}{Department of Mathematics}\par}%
  \vskip 0.18em%
  \tagstructbegin{tag=Title}%
    {\LARGE\bfseries\href{https://gma.math.ufl.edu/exams/analysis-first-year/}{Analysis First-Year Exam}, January 2015, Part 1\par}%
  \tagstructend%
  \vskip 0.35em%
  \tagmcbegin{artifact}%
    \hrule height 1pt%
  \tagmcend%
  \vskip 0.55em%
}
\makeatother
\title{Analysis First-Year Exam, January 2015, Part 1}
\author{}
\date{}
\begin{document}
\maketitle
\thispagestyle{empty}
\begin{examinstructions}
Do exactly 4 problems. Work must be presented in a neat and logical fashion in order to receive credit. Do not leave any gaps. When a theorem is used in a proof it must be precisely stated.
\end{examinstructions}
\begin{examproblems}{0}
\item
Let \((X,d_X)\), \((Y,d_Y)\) be metric spaces. On the Cartesian product \(X\times Y:=\{(x,y):x\in X,y\in Y\}\) define the function 
\[
d_{X\times Y}((x_1,y_1),(x_2,y_2))=d_X(x_1,x_2)+d_Y(y_1,y_2).
\]
 It is known that \(d_{X\times Y}\) is a metric on \(X\times Y\). If~\(X\) and~\(Y\) are complete, must \(X\times Y\) be complete? Prove or give a counterexample.
\item
Let~\(X\) be a metric space and suppose it has the following property: whenever \(\mathcal C\) is a collection of closed subsets of~\(X\), and \(\bigcap_{C\in F}C\) is nonempty for every finite \(F\subset\mathcal C\), then \(\bigcap_{C\in\mathcal C}C\) is nonempty. Prove that~\(X\) is compact.
\item
Let~\(X\) be a metric space and \(A,B\) connected subsets of~\(X\). Prove that if \(A\cap B\) is nonempty, then \(A\cup B\) is connected.
\item
This problem has two parts.
\begin{examsubparts}
\item[\textnormal{a)}]
Define what it means for a function \(f:X\to Y\) to be uniformly continuous on a set \(E\subset X\).
\item[\textnormal{b)}]
Prove that if \(f:X\to Y\) is continuous and~\(X\) is compact, then~\(f\) is uniformly continuous on~\(X\).
\end{examsubparts}
\item
Let \(f:X\to Y\) be a continuous bijection. Prove that if~\(X\) is compact, then \(f^{-1}\) is continuous. Give an example to show that the conclusion can fail if~\(X\) is not assumed compact.
\item
Let \(f:(a,b)\to\mathbb R\) be differentiable and suppose that \(f'\) is monotonically increasing. Prove that \(f'\) is continuous.
\end{examproblems}
\end{document}
